% =====================================================================================
%  Results for the two scoring formulations. Every number here is written by
%  code/formulations/fit_bpi2014.py into code/results_formulations/.
% =====================================================================================

\section{Results: the scoring formulations}
\label{sec:formulation-results}

Section~\ref{sec:results-predictive} measured a constraint. This section measures the two
formulations built under it. The rule for both is the same and is stated once: every constant
is estimated on the earliest \SI{80}{\percent} of the timeline by incident open time, and every
number reported as an evaluation is computed on the remaining \SI{20}{\percent}, which the
fitting never saw. The working set is the \num{39449} incidents and \num{87} groups of
Section~\ref{sec:two-sets}, of which \num{31559} form the training period.

\subsection{What is fitted, and on how much}

Table~\ref{tab:fitted} lists every estimated constant in F1 and F2 with its interval and its
sample size. Nothing in either formulation is a number chosen because it looked reasonable.

\begin{table}[htbp]
\centering
\caption{Every fitted constant in the two formulations. Intervals are \SI{95}{\percent}. All
estimation uses training-period cases only.}
\label{tab:fitted}
\small
%% GENERATED by code/figures/make_tables.py. Do not edit by hand.
%% The caption, label and float live in the manuscript; only the numbers are here.
\begin{tabular}{lL{5.1cm}rl}
\toprule
Symbol & Quantity & Estimate & Interval / support \\
\midrule
$\hat{\lambda}$ & Freshness decay rate, per hour & 0.0174 & [0.0161, 0.0188], $n=\num{31394}$ \\
 & \quad implied half-life (hours) & 39.9 & derived from $\hat{\lambda}$ \\
$\hat{k}$ & Urgency ramp steepness & 0.0228 & $n=\num{31394}$ \\
$\hat{\tau}_0$ & Urgency ramp midpoint, hours & $-12.03$ & $n=\num{31394}$ \\
$\varepsilon$ & Floor applied before any $p\le0$ aggregation & $10^{-3}$ & fixed, reported as a parameter \\
$\alpha$ & Target-encoding smoothing prior & 20 & fixed \\
$\hat{\theta}$ & F2 escalation cut point & 0.5250 & tuned on train, $F_1=\num{0.5134}$ \\
\midrule
 & Cosine-expertise cells evaluated & \num{686430} & $\num{7890}\times\num{87}$ \\
 & \quad of which exactly zero & \num{206214} & \SI{30.04}{\percent} \\
 & Delay rows above the \num{168}\,h cap & \num{171} & \SI{0.4}{\percent} of \num{39449}, a modelling choice \\
\bottomrule
\end{tabular}


\end{table}

Two of these deserve comment rather than a row.

The freshness decay was the first thing to break, and how it broke is worth recording,
because the first explanation we wrote for it was wrong. Fitted across the assignment delay
column as it was then loaded, the maximum-likelihood estimate of $\lambda$ came back exactly
zero with an undefined confidence interval. The delay column showed a median of \num{16.2} hours
against a 90th percentile of \num{5684} hours, and we recorded that as a corrupted upper tail in
the log's first-assignment timestamp.

It was not the log. The loader was subtracting a month-first incident timestamp from a
day-first activity timestamp, so every delay carried a date-convention error of up to eleven
months, and the tail was ours. Parsed correctly the same column has a median of \num{0.49} hours, a
90th percentile of \num{16.57} and a 99th of \num{92.52}. The \num{168}-hour cap now excludes
\num{171} rows of \num{39449}, \SI{0.4}{\percent}, rather than a third of the variable.

The cap therefore survives as a modelling choice and not as a data repair, and it is reported
as one. A wait of more than a week to first assignment is a different operational regime from
the speed-to-lead window $s_5$ is meant to score, so those cases are held out of the fit. The
consequence is measured rather than waved past: the uncapped control fit is now perfectly
well formed at $\hat{\lambda} = \num{0.0097}\,{\rm h}^{-1}$,
interval $[\num{0.0088}, \num{0.0107}]$, so capping
does not rescue a degenerate fit, it nearly doubles the estimated decay rate to
$\num{0.0174}$. Both are reported. Reading $s_5$ off the capped fit assumes the
organisation cares about the first week, which is a statement about the process and not about
the data.

The urgency midpoint $\hat{\tau}_0 = \num{-12.03}$ hours is negative, which means the fitted
logistic ramp is centred before the observed range begins and $s_6$ is therefore monotonically
decreasing across the whole range rather than sigmoidal within it. The shape is real and it is
reported as fitted. It says that on this log urgency has no inflection point inside the
observable window: cases do not become suddenly urgent at some threshold, they become steadily
less recoverable from the moment they arrive.

\subsection{The criterion decision matrix}
\label{sec:decision-matrix}

Table~\ref{tab:criteria} gives the spread of each criterion over the evaluation period. It is
the input to CRITIC, and it explains most of what the weighting analysis then finds.

\begin{table}[htbp]
\centering
\caption{Criterion statistics over the evaluation-period decision matrix. Conflict is the CRITIC
term $\sum_\ell (1-\rho_{k\ell})$. Exact zeros matter because they are what force the floor in
Equation~\ref{eq:powermean}.}
\label{tab:criteria}
\small
%% GENERATED by code/figures/make_tables.py. Do not edit by hand.
%% The caption, label and float live in the manuscript; only the numbers are here.
\begin{tabular}{llrrrrr}
\toprule
 & Criterion & Mean & SD & Min & Exact zeros & Conflict \\
\midrule
$s_1$  & Outcome propensity       & 0.5036 & 0.2958 & 0.0003 & 0 & 5.397 \\
$s_2$  & Service level survival   & 0.5629 & 0.2904 & 0.0000 & \num{16352} & 5.349 \\
$s_3$  & Expertise match          & 0.3608 & 0.4225 & 0.0000 & \num{109012} & 5.996 \\
$s_4$  & Workload headroom        & 0.3301 & $\num{9.8e-06}$ & 0.3301 & 0 & 6.063 \\
$s_5$  & Freshness                & 0.9470 & 0.1289 & 0.0600 & 0 & 6.529 \\
$s_6$  & Urgency                  & 0.4169 & 0.0659 & 0.3676 & 0 & 6.960 \\
$s_7$  & Importance               & 0.1984 & 0.1690 & 0.0010 & 0 & 5.576 \\
\bottomrule
\end{tabular}


\end{table}

\begin{figure}[htbp]
\centering
\includegraphics[width=0.94\textwidth]{fig06_criterion_spread_and_zeros.pdf}
\caption{Criterion distributions and the exact-zero counts that drive the degeneracy in
Equation~\ref{eq:powermean}.}
\label{fig:criteria}
\end{figure}

Three readings, one of them a warning about our own method.

Expertise is bimodal, not continuous. Its standard deviation of \num{0.4225} is the largest in
the table and larger than its mean, which is the signature of a variable that is mostly at zero
or mostly near one. \num{109012} of \num{348000} cells in this decision matrix are exactly
zero because the group has no training history in that incident's category, and the same
measurement over the full evaluation matrix in Table~\ref{tab:fitted} agrees to within a point.
This is the fact that forces $\varepsilon = 10^{-3}$: without it, every $p \le 0$ aggregation would
return zero for roughly a third of all candidate pairs, and the score would silently become a
filter rather than a ranking.

Workload headroom is almost constant at the batch level, standard deviation
$\num{9.8e-06}$. Within a single arrival batch of 40 incidents against \num{87} groups, no assignment
moves Jain's index far. The criterion does real work when accumulated across the whole
optimisation horizon, which is where the objective $f_3$ measures it, but as a per-pair score it
carries almost no discriminating information, and any weighting scheme that reads spread should
be expected to down-weight it.

Freshness sits near its ceiling, mean \num{0.9470}, because most cases in this log are assigned
quickly. A criterion that is near-saturated for most of the data can still be decisive for the
tail, and the exponential form makes it so, but its \emph{average} contribution is small.

\subsection{Weighting, and whether it matters}
\label{sec:weighting-results}

Table~\ref{tab:weights} gives the three weight vectors. They disagree sharply, which is the
expected result: objective weighting methods are known not to be interchangeable
\citep{mukhametzyanov2021specific,ariyanti2025comparative}, to vary in stability under
perturbation \citep{linh2025evaluating,paradowski2025novel} and to admit rank reversal
\citep{li2026trend}.

\begin{table}[htbp]
\centering
\caption{Weights under the three schemes. Entropy weighting assigns effectively zero to
workload headroom, which is the correct response to a criterion with a standard deviation of
$\num{9.8e-06}$ and an illustration of how differently the two schemes read the same matrix.}
\label{tab:weights}
%% GENERATED by code/figures/make_tables.py. Do not edit by hand.
%% The caption, label and float live in the manuscript; only the numbers are here.
\begin{tabular}{lrrr}
\toprule
Criterion & CRITIC & Entropy & Equal \\
\midrule
s1 outcome     & 0.1628 & 0.1283 & 0.1429 \\
s2 sla survival & 0.1589 & 0.1032 & 0.1429 \\
s3 expertise   & 0.2582 & 0.4687 & 0.1429 \\
s4 workload    & 0.1223 & 0.0000 & 0.1429 \\
s5 freshness   & 0.0913 & 0.0071 & 0.1429 \\
s6 urgency     & 0.0783 & 0.0065 & 0.1429 \\
s7 importance  & 0.1283 & 0.2861 & 0.1429 \\
\bottomrule
\end{tabular}


\end{table}

\begin{figure}[htbp]
\centering
\includegraphics[width=0.94\textwidth]{fig05_weights_and_rank_stability.pdf}
\caption{The three weight vectors and the rank correlations between the orderings they induce.}
\label{fig:weights}
\end{figure}

The disagreement between the vectors is large. Entropy puts \num{0.4687} on expertise and
effectively nothing on workload; CRITIC spreads from \num{0.0783} to \num{0.2582}. A reader is
entitled to ask which one the results depend on, and the honest way to answer is to measure it
rather than to argue for a favourite.

The induced assignment orderings agree closely. Spearman $\rho$ is \num{0.959} between CRITIC
and uniform weighting, \num{0.946} between CRITIC and entropy, and \num{0.856} between entropy
and uniform, all at $p < 10^{-10}$. So the weighting choice is not load-bearing for these
results. That is a stronger statement than a defence of any single scheme, and it is
falsifiable in a way that a defence is not: had the correlations come back at \num{0.4}, the
paper would have had to report that the conclusions depend on a modelling choice with no
principled resolution.

We also state why CRITIC alone would have been unsafe even had it agreed. Two of the seven
criteria, $s_1$ and $s_2$, are model outputs. Their spread reflects how the classifier was
calibrated, and isotonic calibration compresses variance, so a CRITIC weight computed on them
measures a property of our pipeline rather than a property of the decision problem. Publishing
that number alone, labelled data-driven, would have laundered an artifact into a result.

\subsection{The compensation exponent}
\label{sec:psweep}

The exponent $p$ in Equation~\ref{eq:powermean} is what turns a qualitative preference into a
measured one. Table~\ref{tab:psweep} sweeps it.

\begin{table}[htbp]
\centering
\caption{Assignment score against the compensation exponent, evaluation period. The floored
column records whether $\varepsilon$ was engaged, which happens for every $p \le 0$ because
\SI{31.3}{\percent} of expertise cells in this matrix are exactly zero.}
\label{tab:psweep}
%% GENERATED by code/figures/make_tables.py. Do not edit by hand.
%% The caption, label and float live in the manuscript; only the numbers are here.
\begin{tabular}{rrrrr}
\toprule
$p$ & Mean score & Min & Max & Exact zeros \\
\midrule
-4.0 & 0.0949 & 0.0012 & 0.5190 & 0 \\
-2.0 & 0.1135 & 0.0014 & 0.6361 & 0 \\
-1.0 & 0.1336 & 0.0018 & 0.7114 & 0 \\
-0.5 & 0.1565 & 0.0031 & 0.7488 & 0 \\
0.0 & 0.2390 & 0.0116 & 0.7828 & 0 \\
0.2 & 0.2918 & 0.0209 & 0.7982 & 0 \\
0.5 & 0.3631 & 0.0558 & 0.8123 & 0 \\
1.0 & 0.4495 & 0.1387 & 0.8371 & 0 \\
2.0 & 0.5494 & 0.2553 & 0.8743 & 0 \\
4.0 & 0.6659 & 0.3602 & 0.9162 & 0 \\
8.0 & 0.7787 & 0.4480 & 0.9504 & 0 \\
\bottomrule
\end{tabular}


\end{table}

\begin{figure}[htbp]
\centering
\includegraphics[width=0.86\textwidth]{fig04_power_mean_sensitivity.pdf}
\caption{Assignment score against the compensation exponent. The curve is monotone in $p$, as
the power-mean inequality requires, which is a useful implementation check as well as a
property.}
\label{fig:psweep}
\end{figure}

The mean score rises monotonically from \num{0.0949} at $p=-4$ to \num{0.7787} at $p=8$, which
is exactly what the power-mean inequality requires and therefore doubles as a correctness test
on the implementation. The operational content is in the minimum column. At $p=1$ the worst
candidate pair in the evaluation set still scores \num{0.1387}, because a strong criterion can
compensate a weak one; at $p=-4$ it scores \num{0.0012}, because it cannot. An organisation
that never wants an incident routed to a group with no relevant history should choose a
negative $p$ and accept that the floor $\varepsilon$ then determines what "no history" is worth.
An organisation willing to trade expertise for availability should choose $p \ge 1$. The
formulation does not decide this; it makes the decision explicit and reports what it costs.

\subsection{F2: deriving breach beats predicting it}
\label{sec:hazard-results}

Table~\ref{tab:hazard} reports the service level formulation. The comparison is between two
models of the same event on the same data, one of which is structurally able to observe its own
label and one of which is not.

\begin{table}[htbp]
\centering
\caption{Service level risk. The person-period expansion turns \num{39449} cases into
\num{161720} rows over ten elapsed-time bins. Every AUC is on the held-out evaluation period.}
\label{tab:hazard}
%% GENERATED by code/figures/make_tables.py. Do not edit by hand.
%% The caption, label and float live in the manuscript; only the numbers are here.
\begin{tabular}{lrl}
\toprule
Model & Held-out AUC & What it estimates \\
\midrule
Direct breach classifier & 0.7879 & $P(\text{breach} \mid \mathbf{x})$ at intake \\
Resolution hazard, bin level & 0.7640 & $h(t \mid \mathbf{x})$ per person-period row \\
Breach derived from resolution, at intake & \textbf{0.8373} & $S(\tau\mid\mathbf{x})/S(t\mid\mathbf{x})$ \\
\midrule
Breach hazard with elapsed time (rejected) & 0.9542 & reads its own label \\
\bottomrule
\end{tabular}


\end{table}

The derived quantity is better than the direct classifier by \num{0.0494} AUC. This is the
result we did not expect and it inverts the usual framing. Admissibility is normally described
as a constraint that costs performance, and Section~\ref{sec:cost} measures exactly that cost.
Here, choosing the object that can actually be estimated \emph{improved} estimation, because the
resolution hazard is a genuine random quantity given elapsed time whereas breach, past the
threshold, is not.

The rejected row is included deliberately. A breach hazard given elapsed time and priority
scores \num{0.9542}, which is by a wide margin the best number in this paper and is worthless.
The reason is a property of the label, and we test it without reconstructing the label's own
threshold, since rederiving a formula only proves it gives the same answer twice. Within every
priority class the longest non-breaching case is shorter than the shortest breaching one, so
breach is a step function of handle time: of the \num{16111} cases at or past that cut, the
measured breach rate is exactly \num{1.0000}. The model is reading the definition of its
label rather than predicting an event. It is reported because a paper about claims that fail
verification should show its own worked example of the failure, and because a reviewer
encountering a number of this size elsewhere in this literature now has a specific mechanism
to test for.

This row is also our own smallest cautionary tale. It stood in the manuscript as a hardcoded
constant for several drafts, carried forward from an earlier run of a dataset we had since
replaced, in the one section of the paper whose subject is numbers that nobody re-derived. It
is now measured in the same run as every other row of Table~\ref{tab:hazard}, on the same
split, with the same features, changing only the event definition.

Two checks are run on the derived quantity rather than assumed. Survival $S(t\mid\mathbf{x})$ is
verified monotone non-increasing over the bin sequence, which it must be by construction and
which fails loudly if the hazard is mis-indexed. And the per-bin resolution rate is
\num{0.2430}, comfortably away from both zero and one, so the hazard model is estimating a
balanced event rather than a rare one.

\subsection{Workload, reported two ways}

Over the \num{87} groups of the evaluation period the load vector has Jain index \num{0.3301},
standard deviation \num{516.7} and coefficient of variation \num{1.424}. Both the bounded and
the unbounded measure are reported for every policy in Section~\ref{sec:assignment-results} and
neither is substituted for the other, because the proposed method loses on workload spread and
replacing the metric it loses on after seeing the result would be fitting the metric to the
outcome.
